\chapter{Cứu Sau Giờ Kiểm Tra}
\chapterintro{Khi lớp vừa căng vừa xẹp}{Những tiết sau kiểm tra rất dễ mất nhịp vì học sinh còn dư mệt và dư cảm xúc.}

\section{Xả một câu rồi học tiếp}
\begin{tinhhuong}{Tình huống}
Sau bài kiểm tra, lớp còn nặng đầu, có em tiếc bài, có em chán học tiếp.
\end{tinhhuong}
\begin{goiydungngay}
Cho cả lớp một vòng thật ngắn: mỗi em nói một từ về trạng thái hiện tại.
\begin{itemize}
\item Chỉ 1 từ hoặc 1 cụm ngắn.
\item Xả cảm xúc xong là đóng lại.
\item Sau đó chuyển sang bài mới rất rõ ràng.
\end{itemize}
\end{goiydungngay}
\begin{cauchot}
Xả đúng lượng thì lớp nhẹ đi; xả quá dài thì tiết học lại trôi tiếp.
\end{cauchot}
\ngancach

\section{Đừng mở bài mới bằng giọng cũ}
\begin{tinhhuong}{Tình huống}
Nếu tiếp tục giảng như không có gì xảy ra, lớp thường nghe ít vì vẫn còn dính tâm trạng cũ.
\end{tinhhuong}
\begin{goiydungngay}
Mở bằng câu: \textit{``Xong phần vừa căng rồi. Giờ mình làm phần này theo kiểu nhẹ hơn.''}
\begin{itemize}
\item Báo hiệu đổi trạng thái rõ ràng.
\item Học sinh cảm thấy được dẫn chuyển nhịp.
\item Rất hiệu quả sau kiểm tra hoặc sau hoạt động mệt.
\end{itemize}
\end{goiydungngay}
\begin{cauchot}
Lớp cần được chuyển trạng thái, không phải bị kéo lê qua trạng thái mới.
\end{cauchot}
\ngancach

\section{Một ý vui để tháo mặt nặng}
\begin{tinhhuong}{Tình huống}
Sau kiểm tra, lớp dễ rơi vào im nặng hoặc than thở kéo dài.
\end{tinhhuong}
\begin{goiydungngay}
Hỏi: \textit{``Nếu bài vừa rồi là một bộ phim, tên phim sẽ là gì?''}
\begin{itemize}
\item Câu hỏi vui giúp hạ căng.
\item Không sa đà bình luận đề quá lâu.
\item Cười một nhịp rồi chuyển bài.
\end{itemize}
\end{goiydungngay}
\begin{cauchot}
Cười đúng lúc là một cách cứu năng lượng học tập.
\end{cauchot}
\ngancach

\section{Tách bài kiểm tra khỏi bài mới bằng một câu rõ}
\begin{tinhhuong}{Tình huống}
Lớp vẫn đang mang cảm giác của bài kiểm tra sang tiết kế tiếp.
\end{tinhhuong}
\begin{goiydungngay}
Nói dứt khoát: \textit{``Phần kiểm tra dừng ở đây. Phần học mới bắt đầu từ đây.''}
\begin{itemize}
\item Câu ngắn nhưng rất có tác dụng chuyển trạng thái.
\item Nói xong thì đưa ngay một nhiệm vụ mới thật nhỏ.
\item Không để lớp đứng giữa hai tâm trạng quá lâu.
\end{itemize}
\end{goiydungngay}
\begin{cauchot}
Lớp đi tiếp tốt hơn khi người dạy cắt ranh giới giữa chuyện cũ và việc mới thật rõ.
\end{cauchot}
\ngancach

\section{Cho lớp thắng một câu dễ trước}
\begin{tinhhuong}{Tình huống}
Sau kiểm tra, nhiều em ngại phát biểu vì sợ tiếp tục sai.
\end{tinhhuong}
\begin{goiydungngay}
Mở phần mới bằng một câu rất dễ hoặc một lựa chọn chắc thắng để lớp có cảm giác làm được ngay.
\begin{itemize}
\item Chỉ cần một chiến thắng nhỏ đầu tiên.
\item Sau đó tăng độ khó dần.
\item Không khí học tập sẽ hồi lại nhanh hơn.
\end{itemize}
\end{goiydungngay}
\begin{cauchot}
Muốn lớp đi tiếp sau một lúc căng, hãy cho các em chạm lại cảm giác làm được.
\end{cauchot}
